\section*{ID6813: S³ (Hypersphere) Eigenmodes — Discreteness, Global Interference, and Why S³ S: κ\_n}
\paragraph{Metadata.}
PiID: 5439478221818 | RTEID: RTE:P26456.3920:T5.8929 | UUID: 4954d8bb-1415-5723-831f-fefa3dc890f3

\paragraph{Canonical Statement.}
If the universe is spatially a 3-sphere of radius \textbackslash{}(R\textbackslash{}) (curved, closed), eigenfunctions of the Laplacian are **hyperspherical harmonics** with discrete eigenvalues \textbackslash{}(\textbackslash{}kappa\_n\textbackslash{}). In coordinates where radial coordinate \textbackslash{}(\textbackslash{}chi\textbackslash{}in[0,\textbackslash{}pi]\textbackslash{}) (great-circle angle), eigenmodes separate as \textbackslash{}[ \textbackslash{}Psi\_\{n,\textbackslash{}ell,m\}(\textbackslash{}chi,\textbackslash{}theta,\textbackslash{}phi) = \textbackslash{}Pi\_\{n\textbackslash{}ell\}(\textbackslash{}chi)\textbackslash{}, Y\_\textbackslash{}ell\textasciicircum{}m(\textbackslash{}theta,\textbackslash{}phi), \textbackslash{}] with radial part \textbackslash{}(\textbackslash{}Pi\_\{n\textbackslash{}ell\}(\textbackslash{}chi)\textbackslash{}) proportional to Gegenbauer-associated functions (spherical Bessel → replaced by \textbackslash{}(\textbackslash{}sin((n+1)\textbackslash{}chi)/\textbackslash{}sin\textbackslash{}chi\textbackslash{}) like structures). Eigenvalues are discrete: \textbackslash{}(k\_n\textasciicircum{}2 \textbackslash{}sim (n+1)\textasciicircum{}2 / R\textasciicircum{}2\textbackslash{}). See e.g. the theory of spherical space harmonics.

\paragraph{Canonical Equation.}
\[
\kappa_n
\]

\paragraph{Dictionary.}
Relation: κ\_n

\paragraph{Encyclopedia.}
S³ (Hypersphere) Eigenmodes — Discreteness, Global Interference, and Why S³ S: κ\_n is a symbolic expression, written κ\_n. It introduces a symbol or relation used elsewhere in the framework. It is built from n. Within the Trawin Zero Point registry it serves as one element of the framework's mathematical narrative.
